\clearpage
\SourceLeaf{175}%
\pdfbookmark[0]{Göttingische gelehrta Anzeigen. 1825 April 11.}{notice1825april11}
\begin{center}
\rule{.82\linewidth}{.35pt}\par
.Göttingische gelehrta Anzeigen.\quad 1825 April 11.\par
\rule{.82\linewidth}{.35pt}
\end{center}
Am 5. April legte Herr Hofrat \textsc{Gauss} der Königlichen Gesellschaft eine Vorlesung mit dem Titel vor:
\begin{center}\itshape
Theoria Residuorum Biquadraticorum, Commentatio prima.
\end{center}

Die Theorie der quadratischen Reste bildet bekanntlich einen der interessantesten Teile der höheren Arithmetik. Nach vielfach wiederholten Untersuchungen kann man sie jetzt als vollendet und abgeschlossen betrachten. Nachrichten zu ihrer Geschichte finden sich in diesen Blättern von 1808 Mai 12 und September 19 sowie von 1817 März 10. An der letzten Stelle wurden auch bereits einige vorläufige Angaben zu den Nachforschungen gemacht, die der Verfasser der vorliegenden Abhandlung seit dem Jahr 1805 über die verwandte, ebenso fruchtbare und interessante, aber viel schwierigere Theorie der kubischen und biquadratischen Reste angestellt hatte. Obwohl er schon damals die wesentlichen Grundzüge dieser Theorie kannte, haben ihn bisher andere Arbeiten daran gehindert, etwas davon zu veröffentlichen. Erst jetzt konnte er sich mit der Ausarbeitung eines Teils dieser Untersuchungen befassen. Den Anfang macht die Theorie der biquadratischen Reste, die mit der Theorie der quadratischen Reste näher verwandt ist als die der kubischen. Die vorliegende Abhandlung soll jedoch
\SourceLeaf{176}%
diesen überaus reichhaltigen Gegenstand noch keineswegs ausschöpfen. Die Entwicklung der \emph{allgemeinen} Theorie, die eine ganz eigene Erweiterung des Gebiets der höheren Arithmetik verlangt, bleibt einer künftigen Fortsetzung vorbehalten. In diese erste Abhandlung sind diejenigen Untersuchungen aufgenommen, die sich ohne eine solche Erweiterung vollständig darstellen ließen. Von den Ergebnissen kann diese Besprechung nur einen Teil hervorheben.

Eine ganze Zahl $a$ heißt biquadratischer Rest der ganzen Zahl $p$, wenn es Zahlen der Form $x^4-a$ gibt, die durch $p$ teilbar sind; biquadratischer Nichtrest dagegen, wenn keine Zahlen dieser Form durch $p$ teilbar sein können. Offenbar sind alle biquadratischen Reste von $p$ zugleich quadratische Reste derselben Zahl und damit alle quadratischen Nichtreste auch biquadratische Nichtreste. Nicht alle quadratischen Reste sind jedoch zugleich biquadratische Reste. Es genügt, die Untersuchungen auf den Fall zu beschränken, dass $p$ eine Primzahl der Form $4n+1$ und $a$ nicht durch $p$ teilbar ist, da sich alle anderen Fälle leicht darauf zurückführen lassen.

Die Untersuchungen dieses Gegenstands gliedern sich in zwei Teile, je nachdem, ob $p$ oder $a$ als gegeben betrachtet wird. Der erste ist viel weniger schwierig als der zweite und im Vergleich zu ihm als ganz elementar anzusehen. Alles Wesentliche, was darüber zu sagen ist, enthält die Abhandlung vollständig.

Aus dem zweiten Teil werden hier dagegen erst einige besondere Fälle behandelt, die sich ohne allzu große Vorbereitungen erledigen ließen und der künftig zu gebenden allgemeinen Theorie vorarbeiten können. Es sind die Fälle, in denen $a=-1$ und $a=\pm2$ gesetzt wird. Der erste bereitet keinerlei Schwierigkeit: Schon in den \emph{Disquisitiones Arithmeticae} wurde gezeigt, dass $-1$ biquadratischer Rest von $p$ ist, wenn $p$ die Form $8n+1$ hat, dagegen nur quadratischer Rest und biquadratischer Nichtrest von $p$, wenn $p$ die Form $8n+5$ hat. Ganz anders verhält es sich mit dem Fall $a=\pm2$. Zwar ist seit Langem bekannt, dass $+2$ und $-2$ quadratische und damit auch biquadratische Nichtreste von $p$ sind, wenn $p$ die Form $8n+5$ hat, und wenigstens quadratische Reste, wenn $p$ die Form $8n+1$ hat. Ebenso ist bekannt, dass bei dieser Form von $p$ entweder $+2$ und $-2$ zugleich biquadratische Reste oder zugleich biquadratische Nichtreste sind:
\SourceLeaf{177}%
Die Entscheidung, welcher dieser beiden Fälle eintritt, erfordert jedoch eine Untersuchung viel höherer Art. Dafür entwickelt die Abhandlung zwei verschiedene Kriterien.

Das erste Kriterium hängt mit der Zerlegung der Zahl $p$ in ein Quadrat und das Doppelte eines Quadrats zusammen. Bekanntlich ist diese Zerlegung (da, wie schon bemerkt, $p$ als Primzahl vorausgesetzt wird) stets und nur auf eine Weise möglich. Setzt man $p=gg+2hh$, so ist $\pm2$ biquadratischer Rest von $p$, wenn $g$ die Form $8n+1$ oder $8n+7$ hat, dagegen biquadratischer Nichtrest, wenn $g$ die Form $8n+3$ oder $8n+5$ hat.

Das zweite Kriterium hängt mit der Zerlegung der Zahl $p$ in zwei Quadrate zusammen, die bekanntlich ebenfalls stets und nur auf eine Weise möglich ist. Setzt man $p=ee+ff$ und nimmt an, dass $ee$ das ungerade und $ff$ das gerade Quadrat bezeichnet, so folgt schon aus der vorausgesetzten Form $p=8n+1$, dass auch $\frac12f$ eine gerade Zahl ist. Also hat $f$ entweder die Form $8m$ oder die Form $8m+4$: Im ersten Fall ist $\pm2$ biquadratischer Rest, im anderen biquadratischer Nichtrest von $p$.

Wir deuten hier nur die Bemerkung an, zu der die höhere Arithmetik so oft Anlass gibt: Nicht so sehr die Schönheit und Einfachheit der Sätze selbst, sondern vielmehr die Schwierigkeit ihrer Beweise macht sie besonders bemerkenswert. Sobald man einen Zusammenhang zwischen dem Verhalten der Zahl $\pm2$ und den beiden angeführten Zerlegungen der Zahl $p$ vermutet, lässt sich dieser Zusammenhang anhand von Zahlenbeispielen äußerst leicht entdecken. Doch schon beim ersten Kriterium ist der Beweis nicht ganz leicht; beim zweiten liegt er viel tiefer verborgen und ist eng mit anderen subtilen Hilfsuntersuchungen verbunden, die ihrerseits zu einer bemerkenswerten Erweiterung der Theorie der Kreisteilung führen. Gerade diese wunderbare Verknüpfung der Wahrheiten ist es. die, wie schon oft bemerkt wurde, der höheren Arithmetik einen so eigenen Reiz verleiht. Diese Beweise selbst lassen sich hier natürlich nicht auszugsweise wiedergeben; man muss sie in der Abhandlung selbst nachlesen. \Needspace{7\baselineskip}Ein paar andere neue arithmetische Sätze, die ebenfalls eng mit dem Beweis des zweiten Kriteriums verbunden sind, verdienen aber wegen ihrer Einfachheit hier noch besonders hervorgehoben zu werden.

Ist $p$ eine Primzahl der Form $4k+1$ und setzt man diese Zahl $=ee+ff$,
\SourceLeaf{178}%
sodass $ee$ das ungerade und $ff$ das gerade Quadrat bezeichnet, und setzt man ferner
\[
1\pdot2\pdot3\mathinner{.\,.\,.\,.\,.\,.}\ k=q
\]
\[
(k+1)(k+2)(k+3)\dotsfive2k=r
\]
so ist stets $\pm e$ der kleinste Rest, der bei Division von $\frac r{2q}$ durch $p$ entsteht, und $\pm f$ der kleinste Rest, den man bei Division von $\frac12rr$ durch $p$ erhält (der kleinste Rest ist dabei stets zwischen den Grenzen $-\frac12p$ und $+\frac12p$ zu nehmen). Die Zahl $\frac r{2q}$, die für $p=5$ den Wert $1$ hat, kann man für größere Werte von $p$ auch in die folgende Form bringen:
\[
\frac{6\pdot10\pdot14\pdot18\dotsfive(p-3)}{2\pdot3\pdot4\pdot5\dotsfive k}
\]
Es ist sehr bemerkenswert, dass sich die Zahl $p$ so auf ganz direktem Weg in zwei Quadrate zerlegen lässt. Fast noch bemerkenswerter ist aber ein dabei auftretender Nebenumstand: Dieses Verfahren liefert die Wurzel des ungeraden Quadrats, $e$, stets mit positivem Vorzeichen, wenn $e$, positiv genommen, die Form $4m+1$ hat, und mit negativem, wenn $e$, positiv genommen, die Form $4m+3$ hat. Für das Vorzeichen, mit dem die Wurzel des geraden Quadrats, $f$, aus dieser Rechnung hervorgeht, hat sich dagegen bisher überhaupt keine allgemeine Regel finden lassen, weder a priori noch anhand von Zahlenbeispielen. Am Schluss der Abhandlung empfiehlt der Verfasser daher diesen Gegenstand den Freunden der höheren Arithmetik zur weiteren Untersuchung, überzeugt, dass ihr Gelingen zugleich eine ergiebige Quelle neuer Erweiterungen dieses schönen Teils der Mathematik erschließen werde.
\BookRule{.28\linewidth}
